\documentclass[12pt]{article}
\usepackage[utf8]{inputenc}
\usepackage[T1]{fontenc}
\usepackage{hyperref}
\title{Seze Language Resources for K-12 Curriculum Development}
\author{Compiled by Claude Sonnet 5.5 (claude-sonnet-5-5) for Andr\'as Kornai}
\date{}
\begin{document}
\maketitle

\section{Introduction}
Seze (ISO 639-3 sze; also Sezo) is an Omotic (Afro-Asiatic) language of the Mao branch, spoken in western Ethiopia near Begi, just north of the Hozo-speaking area. Seze and Hozo are often grouped as ``Begi Mao''. Speaker estimates differ widely: about 3,000 in older sources (Bender 2000) against about 13,000 in a 2015 figure, and Joshua Project figures are higher still. Verbix lists it as threatened and unwritten. Contact with Oromo has reduced Mao-language use in the area.

\section{Primary Educational Resources}
No primer, textbook or school programme was found in search results.

\section{Linguistic and Reference Materials}
The key classification study found is Küspert (2015) on the Mao and Komo languages of the Begi-Tongo area: \url{https://glottolog.org/resource/reference/id/550165}. No dictionary or full grammar was found in these searches. Overview: \url{https://en.wikipedia.org/wiki/Seze_language}. The Endangered Languages Project bibliography: \url{https://endangeredlanguages.com/lang/4607/bibliography}.

\section{Religious and Community Resources}
Joshua Project has a language profile: \url{https://joshuaproject.net/languages/sze}. I did not verify whether any scripture translation exists.

\section{Digital and Audio Resources}
No audio, video or corpus resources were found. Verbix has a brief entry: \url{https://docs.verbix.com/EndangeredLanguages/Sezo}.

\section{Curriculum Development Context}
Seze appears to be unwritten and under-documented. A curriculum would first need an orthography and community-led documentation. Related Mao languages such as Hozo may offer comparative help, but their materials were not searched here.
\end{document}